\documentclass[11pt,a4paper]{article}
\usepackage[margin=23mm]{geometry}
\usepackage{fontspec,xeCJK,amsmath,amssymb,booktabs,hyperref,fancyhdr}
\setmainfont[Path=/home/qiuyong/fonts/,BoldFont=ARIALBD.TTF,ItalicFont=ARIALI.TTF]{arial.ttf}
\setCJKmainfont[Path=/home/qiuyong/fonts/]{1611458310630572.ttf}
\hypersetup{colorlinks=true,urlcolor=blue,linkcolor=blue}
\pagestyle{fancy}\fancyhf{}\fancyhead[L]{OFF23 $\to$ POLAR：重加权推导与方差审计}\fancyfoot[C]{\thepage}
\setlength{\headheight}{15pt}
\setlength{\parskip}{5pt}
\setlength{\emergencystretch}{3em}
\newcommand{\E}{\mathbb E}\newcommand{\Var}{\operatorname{Var}}
\newcommand{\ESS}{\operatorname{ESS}}
\expandafter\def\csname N\endcsname{501}
\expandafter\def\csname TemperatureK\endcsname{300}
\expandafter\def\csname KbtEv\endcsname{0.0258519998}
\expandafter\def\csname DeltaUMeanEv\endcsname{480.863838}
\expandafter\def\csname DeltaUStdEv\endcsname{0.355765091}
\expandafter\def\csname SigmaBeta\endcsname{13.7616082}
\expandafter\def\csname ShiftedRawMean\endcsname{0.0033490793}
\expandafter\def\csname ShiftedRawVarianceDdof0\endcsname{0.0023813603}
\expandafter\def\csname NormalizedWeightVarianceDdof0\endcsname{0.000845860599}
\expandafter\def\csname MeanOneWeightVarianceDdof0\endcsname{212.311856}
\expandafter\def\csname CvSquared\endcsname{212.311856}
\expandafter\def\csname Ess\endcsname{2.34867395}
\expandafter\def\csname EssFraction\endcsname{0.00468797196}
\expandafter\def\csname MaxWeight\endcsname{0.595987077}
\expandafter\def\csname Top3Mass\endcsname{0.932182756}
\expandafter\def\csname LogRawMean\endcsname{-18567.9303}
\expandafter\def\csname LogRawVariance\endcsname{-37130.5026}
\expandafter\def\csname GaussianHeuristicEssFraction\endcsname{5.65592482e-83}

\begin{document}
\begin{center}{\LARGE 重加权估计器：公式推导与数据资格}\\[8pt]
{\large 统计力学推导、501 帧实测结果及 $g(r),c_1,c_2$ 的适用边界}\\[6pt]
2026-10-28 扩展版 · 原 PDF 保留为历史版本\end{center}
\section{先给结论与数据来源}
复算对象是仓库上传的 501 帧水轨迹能量：384 个水分子（1152 原子），生产段 500 ps，每 1 ps 一帧，名义温度 300 K。源势为 MACE-OFF23，目标势为 MACE-POLAR。输入为
\texttt{validation/off23\_500ps\_water1152\_2026-10-23/frames\_501\_total\_energies.csv}。
本报告重算已有能量，不包含新的势函数单点计算，也不认证模型或轨迹的物理精度。

\begin{center}\begin{tabular}{lr}\toprule
量（全部方差采用分母 $n$） & 实测值\\\midrule
帧数 $n$ & \csname N\endcsname\\
$k_BT$ / eV & \csname KbtEv\endcsname\\
$\operatorname{sd}(\Delta U)$ / eV & \csname DeltaUStdEv\endcsname\\
$\beta\operatorname{sd}(\Delta U)$ & \csname SigmaBeta\endcsname\\
相对方差 $\widehat{\mathrm{CV}}^2$ & \csname CvSquared\endcsname\\
归一化权重方差 $v_w$ & \csname NormalizedWeightVarianceDdof0\endcsname\\
均值为 1 的权重方差 $v_R$ & \csname MeanOneWeightVarianceDdof0\endcsname\\
Kish $\widehat{\ESS}$ & \csname Ess\endcsname\\
$\widehat{\ESS}/n$ & \csname EssFraction\endcsname\\
最大单帧权重 & \csname MaxWeight\endcsname\\
前三帧权重之和 & \csname Top3Mass\endcsname\\\bottomrule
\end{tabular}\end{center}
\textbf{推荐回答“方差多大”时报告相对方差 $\widehat{\mathrm{CV}}^2$，并附 ESS。}
原始指数权重的绝对方差随任意能量常数改变；只有给出权重定义才有可比性。当前约 93.2\% 的权重落在三帧上，不能把 501 帧视为 501 个同等信息量的目标系综样本。Kish ESS 是权重集中度诊断，不等同于任意性质的精确独立样本数。
\textbf{数值实现未验收：}实际 MD 日志与 PyTorch OFF 重算的差值，去常数后标准差约
0.030 eV，即 300 K 下约 $1.16k_BT$。这足以影响指数权重。
当前数值是给定 CSV 的复算，尚不能把低 ESS 全部归因于两个势的物理差异。
需同帧、同模型、同精度、同边界条件核实能量和力；不能据此直接认定是 FP32 的全部影响。
\newpage
\section{从配分函数逐步得到重加权公式}
设两种势在相同 $N,V,T$、相同构型空间及测度下定义；动量积分相同，可消去。$x$ 包含所有原子坐标，$L,H$ 表示源势与目标势：
\begin{align}
 Z_L&=\int e^{-\beta U_L(x)}\,dx,&p_L(x)&=Z_L^{-1}e^{-\beta U_L(x)},\\
 Z_H&=\int e^{-\beta U_H(x)}\,dx,&\Delta U(x)&=U_H(x)-U_L(x).
\end{align}
将 $e^{-\beta U_H}=e^{-\beta U_L}e^{-\beta\Delta U}$ 代入，得
\begin{align}
q(x)&=e^{-\beta\Delta U(x)}\quad\text{（未归一化比值；整体常数在自归一化中抵消）},\\
\E_L[q]&=\frac1{Z_L}\int e^{-\beta U_H(x)}dx=\frac{Z_H}{Z_L},\\
\E_H[A]&=\frac{\int A(x)e^{-\beta U_L(x)}q(x)dx}{Z_H}
 =\frac{\E_L[qA]}{\E_L[q]}.
\end{align}
这是一条总体恒等式。对源系综抽到的 $n$ 帧，采用自归一化重要性采样：
\begin{equation}
 \hat w_i=\frac{q_i}{\sum_jq_j},\qquad \hat A_H=\sum_i\hat w_i A(x_i).
\end{equation}
有限样本的比值估计一般有偏；总体恒等式并不保证有限轨迹的无偏性或精度。目标有重要构型而源轨迹未访问时，重加权不能凭空产生这些构型。

\subsection{逐帧修正什么才合理？}
每帧采用自己的真实 $\Delta U_i$ 正是上述公式要求的。若给能量差统一加常数 $C$，
\begin{equation}
q'_i=e^{-\beta C}q_i,\qquad \hat w'_i=\hat w_i.
\end{equation}
因此去掉整体参考能零点不会改善 ESS。反之，逐帧随意减去不同 $C_i$ 会把权重变成 $q'_i=q_i e^{\beta C_i}$，相当于改变目标分布；不能以“对齐能量”为理由消掉真实构型依赖的差值。

例外是有明确证据的计算误差修复（错误单位、原子顺序、能量累加精度等），或用经验证的模型预测真实差值。不能将真实物理差值与数值误差一起任意扣除。当前 OFF 重算能量与 MD 日志间存在数值差异，需同帧同精度复核，不能据此直接宣称全部低 ESS 来自物理差异。
\newpage
\section{四种方差的定义及 ESS 的精确代数关系}
对任意非负 $q_i$，令
\begin{equation}
\bar q=\frac1n\sum_iq_i,\quad v_q=\frac1n\sum_i(q_i-\bar q)^2,
\quad \widehat{\mathrm{CV}}^2=\frac{v_q}{\bar q^2}.
\end{equation}
展开平方可得
\begin{equation}
1+\widehat{\mathrm{CV}}^2
=\frac{n^{-1}\sum_iq_i^2}{(n^{-1}\sum_iq_i)^2}
=\frac{n\sum_iq_i^2}{(\sum_iq_i)^2}.
\end{equation}
Kish 权重集中度定义为
\begin{equation}\boxed{
\widehat{\ESS}=\frac{(\sum_iq_i)^2}{\sum_iq_i^2}
=\frac1{\sum_i\hat w_i^2}
=\frac{n}{1+\widehat{\mathrm{CV}}^2}.}
\end{equation}
这里是\textbf{有限样本的精确代数恒等式}，与帧是否独立无关；把它解释成误差等价独立样本量才需要额外假设。若使用 Bessel 修正的分母 $n-1$，不可直接套同一恒等式。

定义 $R_i=n\hat w_i=q_i/\bar q$，其经验均值严格为 1：
\begin{align}
 v_R&=\frac1n\sum_i(R_i-1)^2=\widehat{\mathrm{CV}}^2
 =\frac n{\widehat{\ESS}}-1,\\
 v_w&=\frac1n\sum_i(\hat w_i-1/n)^2
 =\frac1{n\widehat{\ESS}}-\frac1{n^2}=\frac{v_R}{n^2}.
\end{align}
所以 $v_w$ 数值小不代表重叠好：它被 $n^2$ 缩放了。对比不同帧数时，应优先看 $v_R$ 或 $\ESS/n$。

原始 $q_i$ 在当前约 480 eV 的能量差下会下溢。稳定计算采用
\begin{equation}
a_i=e^{-\beta(\Delta U_i-\Delta U_{\min})},\quad
\hat w_i=a_i/\sum_j a_j.
\end{equation}
此标度下 $\bar a=\csname ShiftedRawMean\endcsname$，$v_a=\csname ShiftedRawVarianceDdof0\endcsname$。
原始量可用对数报告：
\begin{align}
\log\bar q&=-\beta\Delta U_{\min}+\log\bar a
 =\csname LogRawMean\endcsname,\\
\log v_q&=-2\beta\Delta U_{\min}+\log v_a
 =\csname LogRawVariance\endcsname.
\end{align}
对数为自然对数；绝对值极小反映参考能零点，并不表明权重波动小。
\newpage
\section{总体方差、重叠与能量差波动}
真正的密度比为 $r(x)=q(x)/\E_Lq=p_H(x)/p_L(x)$。其总体均值为 1，故
\begin{align}
\Var_L(r)&=\E_L(r^2)-1
 =\int\frac{p_H(x)^2}{p_L(x)}dx-1\\
 &=\chi^2(p_H\Vert p_L).
\end{align}
这给出方差与系综不重叠的直接关系。假设独立同分布、所需矩有限且大数律适用，则
\begin{equation}
\frac{\widehat{\ESS}}n\longrightarrow
\frac1{1+\Var_L(r)}.
\end{equation}
但 $R_i=q_i/\bar q$ 使用经验分母，不是真正已知的 $r(x_i)$；本报告的 212 左右是\textbf{经验相对方差}，不能据此证明总体 $\chi^2$ 也等于该数。缺失的尾部可能使总体方差更大。

\subsection{高斯能量差近似：为何不能直接当实测 ESS？}
若额外假定 $X=\beta\Delta U\sim\mathcal N(m,s^2)$，由正态矩母函数得
\begin{align}
\E(e^{-X})&=e^{-m+s^2/2},&\E(e^{-2X})&=e^{-2m+2s^2},\\
\mathrm{CV}_q^2&=e^{s^2}-1,&\ESS/n&\approx e^{-s^2}.
\end{align}
当前 $s=13.7616$，这个高斯总体近似给出 $\ESS/n\approx5.66\times10^{-83}$，远小于经验值 0.004688。它不是 501 帧的 ESS 计算结果；极端尾部在有限样本中几乎不可检验，高斯模型本身也未获证明。不能把高斯外推当精确的所需帧数预测。

\subsection{KL 与二阶矩尺度：仅作条件模型比较}
Chatterjee--Diaconis 定理依赖目标对数密度比的尾部集中性；其样本量尺度涉及
$D_1=D_{\rm KL}(p_H\Vert p_L)$。而二阶矩涉及
$D_2=\log(1+\chi^2)$，由 Jensen 不等式可得 $D_1\le D_2$。
因此不能称 KL 判据比方差/ESS 判据更严厉，也不是两项独立实验证据。

若\textbf{额外假定完整源分布的} $X=\beta\Delta U\sim\mathcal N(m,s^2)$，
指数倾斜给出 $D_1=s^2/2$、$D_2=s^2$。
代入经验 $s=13.7616$ 仅在这一模型中有 $e^{D_1}\approx1.33\times10^{41}$、
$e^{D_2}\approx1.77\times10^{82}$。501 帧未确定总体尾部；
\textbf{这些不是实测总体散度或真实所需采样帧数。此前将其写成生产算量结论，现予撤回。}

单帧总能量差的标准差为 0.355765 eV，约 $13.76k_BT$。除以固定原子数后仅为 0.308824 meV/atom，但玻尔兹曼因子必须使用\textbf{整盒总能量差}。将每原子能量差直接放入指数会人为抬高 ESS。$\sqrt{N}$ 波动标度还要求局域贡献近似独立或短程相关，不是任意系统都成立的恒等式。

增加样本在适当条件下可改善估计，不能宣称 ESS 必然永远不增长。当前结果支持的是：直接对这 501 帧重加权高度集中，可靠恢复目标统计性质尚未成立。
\newpage
\section{权重方差如何传到统计性质？}
令 $\mu_H=\E_H A$，$m_1=\E_L q$。对独立样本，将比值在总体均值处一阶展开：
\begin{align}
\hat A_H-\mu_H
&\simeq\frac1{nm_1}\sum_{i=1}^nq_i(A_i-\mu_H),\\
\Var(\hat A_H)&\simeq\frac1{nm_1^2}\E_L[q^2(A-\mu_H)^2].
\end{align}
这说明误差还取决于 $A$ 与权重的相关性。相同 ESS 不保证 $g(r)$、能量或密度有相同精度。

对于相关 MD 帧，令 $Y_i=q_i(A_i-\mu_H)$，在平稳、相关可求和条件下：
\begin{equation}
\Var(\hat A_H)\simeq
\frac{\gamma_Y(0)+2\sum_{k\ge1}\gamma_Y(k)}{nm_1^2},
\quad \gamma_Y(k)=\operatorname{Cov}(Y_i,Y_{i+k}).
\end{equation}
应使用连续时间块重采样，并在每个块自助样本内重新归一化权重，完整重算后续性质。不能简单用任意未加权自相关时间去修正 Kish ESS。当前权重集中于极少帧，即使自助法也不能补出未采到的构型。

\subsection{$g(r)$ 与 $c_2$}
固定粒子数及体积下，先定义每帧同样归一化的配对直方图观测量 $A_r(x)$，再做 $\hat g_H(r)=\sum_i\hat w_i A_r(x_i)$。配对计数、壳层体积、有限粒子数归一化必须一致。此前独立脚本的无序对计数因子错误已另行记录，本报告没有替换旧结果。

对于\textbf{均匀、各向同性、单组分标量密度}，OZ 关系给出
\begin{equation}
h=g-1,\quad S(k)=1+\rho\hat h(k),\quad
\hat c^{(2)}(k)=\frac{\hat h(k)}{1+\rho\hat h(k)}.
\end{equation}
固定密度下的线性误差传播为
\begin{equation}
\delta\hat c^{(2)}\simeq\frac{\delta\hat h}{S^2},\qquad
\Var(\hat c^{(2)})\simeq\frac{\Var(\hat h)}{S^4}.
\end{equation}
小 $S$ 区域会放大误差；密度不确定时还须包含密度误差及交叉协方差。分子位点相关需要相应矩阵理论；O--O RDF 不自动等于质心 RDF。

现有独立空间审计发现轨迹沿 $z$ 有明显分层，时间平均分箱密度 CV 约 0.888；因此当前数据\textbf{不能未经论证直接套均匀体相标量 OZ 得到合格的 $c_2$}。
\newpage
\section{$c_1$、参考能零点及修正势的正确含义}
采用 $c_1=-\beta\,\delta F_{\rm ex}/\delta\rho$ 的约定，平衡 Euler--Lagrange 方程为
\begin{equation}
\log(\rho(x)\Lambda^3)-c_1(x)+\beta V_{\rm ext}(x)-\beta\mu=0.
\end{equation}
因此从密度得绝对 $c_1$ 还需要同一约定下的目标化学势。固定 $N$ 的一条轨迹不直接给出跨粒子数的自由能导数。减去空间常数可研究相对剖面，但不能宣称已确定绝对 $c_1$。

\textbf{旧材料中“能量加常数就让 $c_1$ 平移”的说法必须限定：}
若对所有 $N$ 加同一常数 $C$，$F\to F+C$，泛函导数不变；若每分子参考能加 $a$，$F\to F+a\int\rho dx$，则 $c_1\to c_1-\beta a$，$c_2$ 不变。只有在已明确跨 $N$ 的能量基准后才能讨论这种平移。单一固定 $N$ 的 480 eV 偏移不足以判定化学势修正。

修正势可写作
\begin{equation}
U_{\rm corr}(x)=U_L(x)+f_\theta(x),\qquad f_\theta\approx U_H-U_L.
\end{equation}
在 $U_{\rm corr}$ 下\textbf{重新采样}，残差权重为
\begin{equation}
q_{\rm res}=e^{-\beta[U_H-U_{\rm corr}]}.
\end{equation}
若残差的构型依赖波动显著减小且采样覆盖目标，ESS 有望提高；这需要实测验证。只在旧帧上拟合后改写指数、却不保持目标能量定义，不是正确的校正。保留原始目标权重时，拟合本身也不会改变旧轨迹的 ESS。

若使用 MD，修正力须为 $\Delta\mathbf F=-\nabla f_\theta$，力损失正确符号为
$\|\Delta\mathbf F+\nabla f_\theta\|^2$。MC 可以使用能量差，但提议分布及接受率必须满足目标的详细平衡。

\section{文献对照与接下来怎么验证}
Herzog 等的 MLPT 工作讨论源轨迹到目标理论的重加权，并在重叠不足时使用 MLMC 重新采样，而非随意逐帧抹除差值。Hilpert 与 Kresse 的水研究强调势函数能量误差对热力学重加权的影响。两者支持“代理势应接近目标分布”的路线，但不直接认证本数据的 $c_1/c_2$。

建议顺序：同帧统一精度核对两个势及 MD 引擎能量；确认空间分层和目标系综；训练并独立验证修正势后重新采样；最后做权重、时间相关及 $g/S/c_2$ 的完整不确定度评估。论文提供的方法不能替代这几项数据检查。


\newpage
\section{学生势数据的资格与文献解释修订}
\textbf{撤回学生生产尺寸失败结论。}1152 原子的 501 帧来自 OFF23，
学生势只在这些旧帧上补算能量。因此基于 $U_H-U_S$ 的 1.19/501
仅为该指数序列在 OFF 构型上的集中度，不是学生自采样到 POLAR 的 ESS。
如果 $x\sim p_L$，误用该权重得到的分布正比于
$\exp[-\beta(U_L+U_H-U_S)]$，一般不等于 $p_H$。
应采用 $q_{H/L}=q_{H/S}q_{S/L}$，或在学生势下重新平衡并采样。

96 原子案例的 8.43/19 仅为 0.2 ps 短轨迹诊断；气相和 bulk 池化验证集
未被认证为单一平衡系综，不能由评分均匀程度宣称生产统计可用。
未加权学生采样给出学生系综，不能自动作为 POLAR 参考。

\subsection{三种指标不能混用}
对同一批 OFF23\,$\to$\,POLAR 帧，Kish ESS 为 2.34867；
DiffTRe 式(5)的熵形式 $\exp[-\sum_iw_i\log w_i]$ 为 3.06152；
Herzog 式(4)/(5)的升序半权重指标 $I_w$ 为 0。
它们都反映当前经验权重集中，但数值定义与阈值不同，也不证明总体误差条。

\subsection{MBAR 与水论文的解释}
Shirts--Chodera 原文明确允许未直接采样的目标态 $N_i=0$，
前提是已有采样分布覆盖它。不能说每个待估状态必须直接采样；
同样，给旧帧补算中间势能量不会增加构型覆盖。
其最优与无偏陈述属于适当条件下的渐近性质。

Hilpert--Kresse 正文明确报告 63 个水、5000 构型，292 K 时
MACE ESS=4200、KbP ESS=720；323 K 时 KbP ESS=70。
前述“正文未给出4200”的说法撤回。其成功例接近同一个 RPBE-D3 目标，
不直接认证跨 OFF23/POLAR 的本例。截尾可能增大 ESS，但引入偏差和阈值不确定度。

\subsection{2019 MLPT 式(18)的符号不一致}
若预测误差按原文式(17)为 $E\sim\mathcal N(\Delta,s_\epsilon^2)$，则
\begin{equation}
\frac{\int E e^{-\beta E}P(E\mid\Delta)dE}
{\int e^{-\beta E}P(E\mid\Delta)dE}=\Delta-\beta s_\epsilon^2.
\end{equation}
上传式(18)印为正号，与上述定义的积分不一致。
独立 Simpson 积分取 $\beta=1,\Delta=0,s_\epsilon=1$ 得到 $-1$。
这是印刷假设下的符号核对，不推断作者程序如何实现。
$s_\epsilon$ 是预测残差，不能用真实势间差值波动代替；
一般结构量也没有直接加减 $\beta s_\epsilon^2$ 的普遍修正。

\textbf{修订状态：}已修文档解释与构建入口；未重新评估势函数能量，
未生成新的 MD/MC，也未认证目标 $c_1/c_2$。
详细原文页码与复算见仓库的
\texttt{FORMULA\_CROSSCHECK\_PRIMARY\_PDFS\_2026-10-27.md}。

\newpage
\section{实际观测量：氧--氧对数与质心密度}
本节与当前源程序一一对应，不将候选估计器标为已实现。设 $x$ 为一帧全部坐标，
$r_{ij}(x)$ 为周期盒中两氧原子的最小镜像距离。第 $k$ 个壳层的实际无序对计数为
\begin{equation}
H_k(x)=\sum_{1\le i<j\le N_O}
\mathbf 1\!\left[r_k\le r_{ij}(x)<r_{k+1}\right].
\end{equation}
只有落在 $[r_k,r_{k+1})$ 的氧原子对才贡献 1；$i<j$ 保证每一对只计一次。
壳层体积与无序对的均匀参考值为
\begin{equation}
\Delta V_k=\frac{4\pi}{3}(r_{k+1}^3-r_k^3),\qquad
C_k=\frac{N_O(N_O-1)}{2V}\Delta V_k.
\end{equation}
当前轨迹固定 $N_O=384$、固定体积 $V$，所以实际每帧观测量可写为
\begin{equation}\boxed{
A_k^{OO}(x)=\frac{H_k(x)}{C_k},\qquad
\hat g_{OO,H}(r_k)=\frac{\sum_iq_iA_k^{OO}(x_i)}{\sum_iq_i}.}
\end{equation}
当前代码计算 $\sum_iw_i H_k(x_i)$，并除以相同权重下的参考计数；固定 $N_O,V$ 时正等价于上式。
程序使用 $0$--$8\,\text{\AA}$、80 个壳层，区间宽度 $0.1\,\text{\AA}$。
对变体积、变粒子数或约束变化的其他体系，$C_k$ 不再是常数，不能直接套上述固定盒等价关系。
源码为 \texttt{validation/off23\_500ps\_water1152\_2026-10-23/\allowbreak reweight\_500ps\_water1152.py}。

氧位点 $g_{OO}$ 不是分子质心 $g$。若希望为分子 cDFT 提供一体密度，
对于截面积 $S$、厚度 $\Delta z$ 的 $k$ 区间，可选用
\begin{equation}
A_k^{\rho}(x)=\frac{1}{S\Delta z}\sum_{m=1}^{N_m}
\mathbf 1\!\left[z_k\le Z_m^{\rm COM}(x)<z_{k+1}\right].
\end{equation}
这里 $Z_m^{\rm COM}$ 是第 $m$ 个分子的质心坐标；目标质心密度是 $\E_H A_k^{\rho}$。
现有空间诊断使用氧位置，并非此质心估计器。
实际水轨迹沿 $z$ 明显不均匀，\textbf{上述球壳归一化的氧对曲线是诊断，不能自动视为均匀体相 $g$ 或合格 $c_2$}。

\newpage
\section{分子分母使用不同采样频率}
定义 $q(x)=e^{-\beta[U_H(x)-U_L(x)]}$。若使用 $m$ 个含结构量的帧
$S_A$、$M$ 个含权重的帧 $S_q$，可考虑
\begin{equation}\boxed{
\hat\mu_k^{\rm multi}=\frac{\hat N_k}{\hat D},\quad
\hat N_k=\frac1m\sum_{i\in S_A}q(x_i)A_k(x_i),\quad
\hat D=\frac1M\sum_{j\in S_q}q(x_j).}
\end{equation}
分子中的 $A_k$ 与 $q$ 必须在同一构型计算；两组构型必须代表同一平衡源系综，
并有充分覆盖。可令 $S_A\subseteq S_q$，但不同时间点采样也要求恰当的平稳性与关联分析。
有限样本比值一般仍有偏；若帧数不等，不能将分子的 $q_i$ 直接除以分母原始总和而漏掉 $M/m$。

较密地评估 $q$ 可能改善分母的时间积分，但 $q(x)$ 是构型上的函数，也可能包含慢的几何关联；
其关联时间及方差须由数据测量。固定总轨迹长度下，加密输出并不自动产生同样多独立样本。
评估一个新的 $q$ 需要新构型的 $U_L$ 与 $U_H$，特别是 POLAR 单点的成本必须计入。
当前只存每 1 ps 的构型及两套能量，共 501 帧，不能从这些文件恢复 1 ps 之间真实的 $q$。

令 $N_k=\E_L(qA_k)$、$D=\E_Lq$、$\mu_k=N_k/D$，写成样本波动 $\delta N_k,\delta D$ 后，
一阶误差展开为
\begin{align}
\delta\hat\mu_k&\simeq\frac{\delta\hat N_k-\mu_k\,\delta\hat D}{D},\\
\Var(\hat\mu_k)&\simeq\frac{\Var(\hat N_k)+\mu_k^2\Var(\hat D)
-2\mu_k\operatorname{Cov}(\hat N_k,\hat D)}{D^2}.
\end{align}
同一条轨迹的分子与分母共享帧时，协方差可以抵消或增加误差；多采分母并不保证比同帧比值更优。
应同时检查 $q$、$A_k$ 和 $qA_k$ 的时间相关，并以共同时间块对两组采样联合重采样。
只比较分子、分母各自的误差条不够。

\newpage
\section{分部积分：用目标力估计密度梯度}
以下先针对可微势、正则系综、笛卡尔坐标与可消去积分边界项的情形；
硬墙、刚性约束、有限边界或特殊坐标时需要补充对应边界和测度项。
令 $Z_m$ 为粒子或分子的位置，$F^H_{m,z}$ 为目标势产生的相应总力的 $z$ 分量。
对于截面积 $S$ 的平面平均密度，分布意义下
\begin{equation}
\rho_H(z)=\frac1S\left\langle\sum_m\delta(z-Z_m)\right\rangle_H.
\end{equation}
对 $z$ 求导，并用 $\partial_z\delta(z-Z_m)=-\partial_{Z_m}\delta(z-Z_m)$，
再对构型坐标分部积分，得到
\begin{equation}\boxed{
\frac{d\rho_H(z)}{dz}=
\frac{\beta}{S}\left\langle\sum_m F^H_{m,z}\delta(z-Z_m)\right\rangle_H.}
\end{equation}
其中 $F^H=-\nabla U_H$。如果 $Z_m$ 是分子质心，$F^H_{m,z}$ 应使用该分子受到的总力，
并检查所选刚性或柔性模型的约束与坐标测度。若是氧位点密度，则应采用氧原子的相应总力。
在周期 $z$ 方向，积分常数需用粒子数归一化 $\int\rho_H(z)S\,dz=N$ 确定。

在源轨迹上估计目标力密度时，对含力的每帧采用
\begin{equation}
\left\langle\sum_mF^H_{m,z}\delta(z-Z_m)\right\rangle_H
=\frac{\left\langle q\sum_mF^H_{m,z}\delta(z-Z_m)\right\rangle_L}
{\E_Lq}.
\end{equation}
特别地，不能只把源力 $F^L$ 乘 $q$ 当成目标力：
$F^H=F^L-\nabla\Delta U$，还需要势差的力修正。
对不同区间近似 $\delta$ 并积分后，需要检查粒子数、周期连续性以及可能的负密度等误差。
\textbf{此方案在仓库尚未实现，也没有完成低方差验证。}

力采样可减少窄箱直方图的统计噪声，但不保证在每个体系和每种权重下都更高效；
目标力计算成本、长程力精度与 $qF^H$ 的相关时间都应计入。
可参考 de las Heras 与 Schmidt 的一体密度力采样，以及 Rotenberg 关于密度和径向分布估计的综述。

\newpage
\section{应如何验证估计器及当前数据限制}
\textbf{方案 A：}保持原有同帧直方图比值，作为基线；
\textbf{方案 B：}分别设定分母 $q$ 与分子 $qA$ 的采样间隔；
\textbf{方案 C：}使用目标力密度重建一体密度，或先以可验证的体相体系推导/测试力辅助 RDF。
对 C，先给出与粒子类型、约束、边界一致的完整 RDF 推导，不将密度梯度公式直接当作 RDF 公式。

对各方案在相同状态和总计算预算下记录：重加权源轨迹长度、源势/目标势的能量和力调用次数、
墙钟时间、各距离区间或空间区间的偏差、联合分块误差和对采样间隔的敏感性。
基于有效样本、相关时间和观测量的误差--成本曲线选择方案；
Kish ESS 只是权重集中度，不能代替 $g(r)$ 或密度的误差条。

当前公开能量表记录 $U_L,U_H$ 与同一批构型，但不含逐原子 POLAR 目标力；
因此仅能检查现有同帧结果，并在现有时间分辨率内比较下采样。
对更密的 $q$ 或目标力估计，必须有对应构型并重新执行目标势单点，不能用插值冒充新数据。

上面的目标平均恒等式允许估计 $g$ 或 $\rho$；从 $\rho$ 到绝对 $c_1$ 还需目标化学势及一致参考。
从 $g$ 到 $c_2$ 则要求与均匀性、位点/分子表示和有限尺度相匹配的关联关系。
\textbf{此处提出了可比较的估计器和验证方案，没有声称新方法已计算、已经提高精度或已完成 $c_1/c_2$ 标签认证。}

\newpage
\begin{thebibliography}{9}
\bibitem{herzog} B. Herzog et al., Assessing the Accuracy of Machine Learning Thermodynamic Perturbation Theory: Density Functional Theory and Beyond. JCTC 18, 1382--1394 (2022). \href{https://doi.org/10.1021/acs.jctc.1c01034}{DOI: 10.1021/acs.jctc.1c01034}.
\bibitem{hilpert} T. Hilpert and G. Kresse, Accurate Thermophysical Properties of Water using Machine-Learned Potentials. JCP 164, 194504 (2026). \href{https://doi.org/10.1063/5.0325444}{DOI: 10.1063/5.0325444}; \href{https://arxiv.org/html/2601.21103v1}{开放全文}.
\bibitem{cd} S. Chatterjee and P. Diaconis, The Sample Size Required in Importance Sampling. Ann. Appl. Probab. 28, 1099--1135 (2018). \href{https://doi.org/10.1214/17-AAP1326}{DOI: 10.1214/17-AAP1326}.
\bibitem{force} B. Rotenberg, Use the force! Reduced variance estimators for densities, radial distribution functions and local mobilities in molecular simulations. JCP 153, 150902 (2020). \href{https://doi.org/10.1063/5.0029113}{DOI: 10.1063/5.0029113}.
\bibitem{density} D. de las Heras and M. Schmidt, Better than counting: Density profiles from force sampling. PRL 120, 218001 (2018). \href{https://doi.org/10.1103/PhysRevLett.120.218001}{DOI: 10.1103/PhysRevLett.120.218001}.
\bibitem{ess} V. Elvira, L. Martino and C. P. Robert, Rethinking the Effective Sample Size. International Statistical Review (2022). \href{https://doi.org/10.1111/insr.12500}{DOI: 10.1111/insr.12500}.
\end{thebibliography}
\end{document}
